\documentclass[11pt, a4paper]{article}
\usepackage{graphicx}
\usepackage{booktabs}
\usepackage{amsmath}
\usepackage{url}

\title{\textbf{Metabolic State Clustering in Meningioma Single-Cell RNA-seq Data}}
\author{}
\date{}

\begin{document}

\maketitle

\section*{Abstract}
This study addresses whether transcriptional metabolic states can be used to cluster or separate groups in a single-cell RNA-seq dataset of 58,843 cells from 10 meningioma samples. We analyze 1,984 metabolic genes from Reactome pathway R-HSA-1430728 (release 97) after aggregating duplicate symbols and filtering for prevalence.

\section{Methods}
\subsection{Data Processing}
Raw UMI counts were normalized by per-cell library size to 10,000, then log-transformed using the natural logarithm:
\begin{equation}
\label{eq:log1p}
\log(1+x) = \log(1 + x)
\end{equation}
where $x$ is the normalized count and $\log(1+x)$ is the log1p-transformed value.

\subsection{Gene Set Selection}
Reactome metabolism pathway R-HSA-1430728 (release 97) provided 2,197 genes. After filtering for prevalence ($\ge 20$ cells) and excluding mitochondria, 1,984 genes remained. Duplicate symbols were aggregated by summing raw counts before QC.

\subsection{Clustering}
Leiden clustering was performed with resolution 0.4, selected via stability filtering (mean pairwise ARI $\ge 0.75$) across seeds $[0, 17, 42, 73, 101]$ and resolutions $[0.4, 0.8, 1.2]$. Feature scaling used Scanpy's scale with $zero\_center=False$ and $max\_value=10$.

\subsection{Baseline Comparisons}
HVG-based clustering and 19 size-matched random controls were compared. The metabolic gene set achieved silhouette score 0.229, with 0 of 19 random controls exceeding or equal to this value (add-one rank fraction: 0.05).

\section{Results}
\subsection{Global Clustering Analysis}
Leiden clustering with 15 clusters (sizes: 1,002--8,985 cells) yielded a silhouette score of 0.229. Cluster sizes ranged from 1,002 to 8,985 cells.

\subsection{Confounder Assessment}

NMI values for categorical confounders:
\begin{table}
\toprule
Confounder & NMI \\
\midrule
patient & 0.544 \\
sample & 0.567 \\
source & 0.197 \\
cell type & 0.507 \\
cell subtype & 0.658 \\
cell cycle phase & 0.195 \\
\bottomrule
\end{table}

\subsection{Within-Replicate Sensitivity}
Across 6 patients, the weakest replicate (patient 3, 3,287 cells) showed global ARI 0.106 and within-replicate ARI range 0.106--0.813.

\section{Discussion}
The analysis found candidate partitions but distinct metabolic states remain inconclusive. Associations are descriptive and non-causal. CD8 candidate partitions were observed but distinct metabolic states remain inconclusive.

\section{References}
\begin{enumerate}
\item Choudhury, A. et al. (2022). \textit{Meningioma DNA methylation groups identify biological drivers and therapeutic vulnerabilities}. Nature Genetics, 54(5), 649--659. \url{https://doi.org/10.1038/s41588-022-01061-8}
\end{enumerate}

\end{document}
